\section{Volume estimates uniform in the bare offset}
\label{sec:volume}

The trajectories admitted in Proposition~\ref{prop:admission} have
terminal couplings in a fixed band, but generally do not end at $H$.
We now extend the finite-volume estimates of Section~\ref{sec:trace}
to this entire family.  The argument first matches two trajectories
at their common terminal coupling, then transfers a rectangular
partition-function test from a retuned reference depth to every finer
cutoff.  The resulting bounds will allow local field comparisons on
fixed tori to pass to the plane.

\subsection{Positivity at finite coupling}
\label{sec:finite-beta}

We first record the elementary positivity needed to use component
moments as measures.  All torus partition functions retain the
normalization in \eqref{eq:lattice-model}: normalized area measure at
each site, and each unoriented nearest-neighbor bond counted once.

\begin{lemma}[Finite-coupling moments and normalization]
\label{lem:finite-beta}
For every sufficiently large fixed $\beta$, the periodic plane law
$\mu_\beta$ exists and is translation and $O(3)$ invariant.  Every
component moment
\[
 \E_{\mu_\beta}\prod_{r=1}^n q_{x_r}^{i_r},
 \qquad x_r\in\Z^2,\quad i_r\in\{1,2,3\},
\]
is nonnegative.  If
\[
 C_\beta(x)=\E_{\mu_\beta}[q_0^1q_x^1],\qquad
 \chi_\beta=\sum_{x\in\Z^2}C_\beta(x),\qquad
 \xi_\beta^2=\frac{\sum_x|x|^2C_\beta(x)}{4\chi_\beta},
\]
then both sums converge absolutely and
$0<\chi_\beta,\xi_\beta<\infty$, where $\xi_\beta$ is the positive
square root.  More precisely, $C_\beta(0)=1/3$ and, for each
nearest-neighbor vector $e$,
\begin{equation}
 C_\beta(e)\ge\frac{\beta}{9}e^{-7\beta}>0.
 \label{eq:offset-edge-positive}
\end{equation}
\end{lemma}

\begin{proof}
Proposition~\ref{prop:preliminary} gives the periodic plane limit.
The symmetries of its torus approximants pass to that limit; in
particular each component has mean zero.  Apply
\eqref{eq:prelim-mixing} to $q_0^1,q_x^1$ on arbitrarily large tori
at this fixed $\beta$, and then pass to the plane.  It gives
\[
 |C_\beta(x)|\le
 C(1+\beta+|x|)^p\exp[-c|x|/X(\beta)].
\]
The two sums in the statement therefore converge absolutely.

For positivity, expand each factor
$\exp(\beta\sum_{i=1}^3q_x^iq_y^i)$ in its componentwise power series
on a finite torus.  The series of absolute values is integrable,
since the spins are bounded and the torus has finitely many edges.
Every integrated term, including any prescribed component insertion,
factors into one-site monomial integrals.  Such an integral vanishes
if one coordinate exponent is odd and is nonnegative if all are
even.  All expansion coefficients are nonnegative.  Division by the
positive partition function, followed by the periodic limit, proves
the asserted moment positivity.

Consider one edge $\{0,e\}$ on a torus with sides at least four.
Exactly seven distinct bonds meet its two endpoints.  Delete these
bonds and denote the resulting partition function by $Z^{\rm del}$.
In the expansion of the numerator for $q_0^1q_e^1$, retain the term
$\beta q_0^1q_e^1$ from the chosen edge and the constant terms from
the other six incident bonds, leaving all nonincident expansions
unrestricted.  The isolated endpoint integrals each equal $1/3$,
so this contribution is $\beta Z^{\rm del}/9$.  Termwise
nonnegativity bounds the full numerator from below by this value.
Since the deleted interaction is at most $7\beta$ pointwise,
$Z\le e^{7\beta}Z^{\rm del}$.  This proves
\eqref{eq:offset-edge-positive}, also after passing to the plane.
Internal symmetry gives $C_\beta(0)=1/3$.  Nonnegativity now gives
$\chi_\beta\ge1/3$ and a strictly positive second-moment numerator.
\end{proof}

\subsection{Endpoint comparisons with a varying terminal coupling}

The reference scale at the endpoint remains $H$ throughout.  Thus the
thresholds are always $t_0=H^{-1/2}p_0$, the regular-error cap is
$\delta_0=H^{-2.05}$, and the covering cap is
$w_0=\exp(-p_0^{1/4})$.  These quantities are not replaced by
functions of the precise terminal coupling.

\begin{lemma}[Ordinary endpoint laws throughout the terminal band]
\label{lem:band-endpoints}
Fix a finite covering-cap multiplier and $C_0>0$ before choosing $H$
sufficiently large.  Let a trajectory obtained from the microscopic
Gibbs law by the exact ordinary observations satisfy the admission and
cap conditions of Section~\ref{sec:setup}, with terminal coupling
$b_0\in[H/2,2H]$.  On a standard terminal axis torus of bounded
aspect, with the lower bounds and divisibilities of
Section~\ref{sec:comparison}, its scalar-stripped endpoint density
is an admissible base for Proposition~\ref{prop:integrated-comparison}
and Corollary~\ref{cor:terminal-alignment}, with constants uniform in
the trajectory, its depth, its period, and $b_0$ in this band.
In particular, every endpoint bond obeys
\begin{equation}
 \mu\{|q_x-q_y|>t_0/C_0\}\le e^{-cHt_0^2}
 \label{eq:offset-terminal-alignment}
\end{equation}
for a corresponding $c>0$.
The integrated comparison also permits complex comparison weights
whose support-hit envelope obeys the chosen multiplier of the covering
cap, as in its original statement.
\end{lemma}

\begin{proof}
The definition in Section~\ref{setup:endpoints}, specifically
\eqref{setup:endpoint-law}, already allows a precise coupling in
$[H/2,2H]$.  An actual observed law is positive because its density
is the integral of the positive microscopic density against positive
normalized observations.  The ordinary observations on the two
sides of a retained seam are independent and are exchanged by
reflection.  The verification at the end of
Section~\ref{sec:shooting} therefore supplies link reflection
positivity without any condition $b_0=H$.

The quantitative estimates used in the integrated comparison have
the same uniformity.  In Lemma~\ref{lem:comparison-remainder},
group the kinetic rows with their precise coupling $b_0$.  The omitted
tails cost at most $2H$ times the exponentially small kernel tail,
and the canonical contributions remain bounded by
$C_L(Ht_0^4+t_0^2)$ per anchor.  In
Lemma~\ref{lem:editable-regions}, the lower and upper energy bounds
use respectively $b_0\ge H/2$ and $b_0\le2H$; their energy scale is
still $Ht_0^2=p_0^2$.  The sampling, masks, complete supports, and
inventory constraints have the same reference thresholds and caps.
The chessboard argument uses only the seam reflection positivity and
the even tile counts, both just checked.  These are precisely the
hypotheses of Proposition~\ref{prop:integrated-comparison}.
The disseminated-event proof of
Corollary~\ref{cor:terminal-alignment} uses the same remainder and
partition-function bounds and the lower kinetic bound $b_0\ge H/2$.
It therefore gives \eqref{eq:offset-terminal-alignment} uniformly.
\end{proof}

\begin{proposition}[Endpoint comparison at a common terminal value]
\label{prop:endpoint-comparison}
Fix $L,R,R'$ and choose $H$ sufficiently large as in
Proposition~\ref{prop:admission}.  Consider two admitted ordinary
trajectories of depths $N\ge K\ge1$, with offsets bounded by $R'$,
common reference scales and caps, and the same precise terminal
coupling $b_0$.  Align their last $K$ layers on common compatible
periods.  Write their scalar-stripped endpoint densities as
$e^{X_i}\Xi_i$.  On each common admitted terminal torus of volume
$v$, their common refined activity lists satisfy
\begin{align}
 \|X_1-X_2\|_\infty
    &\le C_Hv\epsilon_K,\notag\\
 \sup_x\sum_{\lambda:x\in P_\lambda}
 e^{As_\lambda}\|k_{1,\lambda}-k_{2,\lambda}\|_\infty
    &\le C_H\epsilon_K,
 \qquad \epsilon_K=L^{-(2-\upsilon)K}.
 \label{eq:offset-endpoint-comparison}
\end{align}
Here $P_\lambda$ and $s_\lambda$ are the complete support and load
of a covering label, and $A$ is the support exponent of the endpoint
class.  If the common torus is an axis torus satisfying
Lemma~\ref{lem:band-endpoints} and $C_Hv\epsilon_K$ is sufficiently
small, its scalar-stripped partition functions satisfy
\begin{equation}
 |\log\widehat Z_1-\log\widehat Z_2|
       \le C_Hv\epsilon_K.
 \label{eq:offset-endpoint-logZ}
\end{equation}
The constants are independent of $N,K$, the periods, and the common
terminal value.
\end{proposition}

\begin{proof}
We verify the uniformity in the common terminal value in the
two-boundary proof of Proposition~\ref{prop:endpoint-matching}.
The normalized discrepancies $u_j$ and precise-coupling differences
$\lambda_j$ are those of \eqref{eq:setup-density-discrepancy}.
The equality of the terminal couplings is exactly $\lambda_0=0$.
The shared ordinary tail gives history error at most
$r_0^{K-j}$, by Proposition~\ref{prop:free-bounds}, with
$r_0=A_0/L^2$.  Thus Theorem~\ref{thm:rg-step} gives, for
$1\le j\le K$,
\[
 u_{j-1}\le q u_j+e_*|\lambda_j|+f_j,
 \qquad
 |\lambda_j|\le a_*|\lambda_{j-1}|+a_*C_Lu_j+a_*f_j,
\]
where
\[
 e_* = \sup_{j\ge1}C_L(\log H_j)^C/H_j=o_H(1),
 \quad a_*=(1-e_*)^{-1},\quad
 f_j\le C_L\operatorname{poly}(H_j)r_0^{K-j}.
\]
Choose $\max(q,r_0)<\rho<L^{-2+\upsilon}$ with strict spare width,
and put $F_K=C_H(1+K)^C$, large enough that
$f_j\le F_K\rho^{K-j}$.  For
\[
 U=\max_{0\le j\le K}\frac{u_j}{\rho^{K-j}},
 \qquad W=\max_{0\le j\le K}
                  \frac{|\lambda_j|}{\rho^{K-j}},
\]
downward iteration for $u$ and upward iteration from $\lambda_0=0$
give
\[
 U\le u_K+\frac{e_*W+F_K}{\rho-q},
 \qquad
 W\le\frac{a_*(C_LU+F_K)}{1-a_*\rho}.
\]
Choose $H$ so that $a_*\rho<1$ and the product of the two
coefficients coupling $U,W$ is less than $1/2$.
The common caps bound $u_K$, so
$u_j+|\lambda_j|\le C_H(1+K)^C\rho^{K-j}$.
This calculation uses no specified value of $b_0$ beyond its band.

At layer zero the conversion in the proof of
Proposition~\ref{prop:endpoint-matching} uses the coefficient norms,
masked regular norms, complete-support load bounds, and free-history
kernel bounds.  These are the common reference norms just used.
It gives \eqref{eq:offset-endpoint-comparison}, absorbing the factor
$(1+K)^C$ in the spare exponential width.  All regular supremum bounds
remain on their prescribed masks; common-list refinements keep their
exact covering corrections.

For the partition functions, set
$\varepsilon=C_Hv\epsilon_K$.  Summing the support-hit bound gives
\[
 \sum_\lambda e^{2s_\lambda}
       \|k_{2,\lambda}-k_{1,\lambda}\|_\infty\le\varepsilon.
\]
By Lemma~\ref{lem:band-endpoints} and
Proposition~\ref{prop:integrated-comparison},
\[
 \left|
 \frac{\int e^{X_1}\Xi_2\dd q}{\widehat Z_1}-1
 \right|\le e^{\varepsilon}-1.
\]
Both actual densities are positive, hence $\Xi_2>0$.
The exponent bound therefore compares $\widehat Z_2$ with
$\int e^{X_1}\Xi_2\dd q$ by factors between
$e^{-\varepsilon}$ and $e^{\varepsilon}$.  For small $\varepsilon$
taking logarithms proves \eqref{eq:offset-endpoint-logZ}.
Each extracted bulk scalar is exactly per volume and independent of
the period, by \eqref{eq:setup-effective-density}; winding
corrections remain in the density.  This fact will allow the scalars
to cancel in the doubling test.
\end{proof}

\subsection{One reference box for every offset}

Recall the rectangular test from \eqref{eq:trace-doubling-definition},
\[
 \Delta_\beta(n,w)=4\log Z_\beta(n,w)-\log Z_\beta(2n,2w).
\]
For later plane comparisons we use
\begin{equation}
 \begin{split}
 \mathcal A&=\{(1,1),(1,25)\},\\
 \mathcal U&=\mathcal A
 \cup\{(u_1/2,u_2):(u_1,u_2)\in\mathcal A\}
 \cup\{(2u_1,u_2/2):(u_1,u_2)\in\mathcal A\}.
 \end{split}
 \label{eq:offset-aspect-inventory}
\end{equation}
The auxiliary aspects are needed when one doubles the time period
and then the spatial period.  The aspect $(1,25)$ will also cover
the inclined microscopic tori.

\begin{lemma}[A reference-box test uniform in the offset]
\label{lem:offset-box-test}
The contraction slack $\upsilon>0$ and then the parameters may be
chosen so that there exist integers $K_0,M_0\ge1$ with
\begin{equation}
 0\le\Delta_\beta(u_1M_0L^N,u_2M_0L^N)\le2^{-10}
 \label{eq:offset-box-test}
\end{equation}
whenever $N\ge K_0$, $\beta=\vartheta_N+s$, $|s|\le R$, and
$(u_1,u_2)\in\mathcal U$.  The integer $M_0$ may be required to
exceed any fixed lower bound and to be divisible by any fixed
integer.  In particular all displayed half-periods and all tile
counts required in Section~\ref{sec:comparison} can be even.
\end{lemma}

\begin{proof}
Let $\eta>0$ be the gain in Proposition~\ref{prop:preliminary}.
Choose the slack small enough to admit a number $z$ with
\begin{equation}
 0<\upsilon/2<z<\min\{1,\eta/(2\pi)\}.
 \label{eq:offset-box-exponents}
\end{equation}
Fix all large-$H$ requirements, including
Lemma~\ref{lem:band-endpoints}, before choosing the box sizes.
Take $m_K$ to be a fixed common integer multiple of a rounded
dyadic integer with
\[
 \log_L m_K=(1-z)K+O(1).
\]
The fixed multiple clears the denominators of $\mathcal U$ and
imposes the divisibilities and lower bounds in the statement.

For each depth-$N$ ordinary run with $|s|\le R$, and each $K\le N$,
Lemma~\ref{lem:retuning} supplies a depth-$K$ ordinary run with
bare coupling $\widetilde\beta_K=\vartheta_K+\widetilde s_K$,
$|\widetilde s_K|\le R'$, and exactly the same terminal coupling.
Uniformly over these retunings,
\[
 \widetilde\beta_K=H+\gamma K+O_{H,L,R'}(\log(2+K)),\qquad
 X(\widetilde\beta_K)
 \le C_H(1+K)^C L^{(2-\eta/(2\pi))K}.
\]
The shorter side of $(u_1,u_2)m_KL^K$ is comparable to
$L^{(2-z)K}$.  Its ratio to this preliminary upper length grows
exponentially, by $z<\eta/(2\pi)$.  The even periods thus satisfy
all hypotheses of \eqref{eq:prelim-doubling}, whose exponentially
small factor dominates its polynomial prefactor.  Consequently
\[
 \sup_{\substack{N\ge K,\ |s|\le R\\(u_1,u_2)\in\mathcal U}}
 \Delta_{\widetilde\beta_K}(u_1m_KL^K,u_2m_KL^K)
       \longrightarrow0.
\]

Run both cutoffs to their common terminal rectangle
$(u_1,u_2)m_K$, of volume $v=u_1u_2m_K^2$, and do the same on
the doubled rectangle.  These periods are admitted: the imposed
terminal lower bound and the growth of the periods through earlier
layers give the conditions in Section~\ref{setup:scales}.
Proposition~\ref{prop:endpoint-comparison} applies, because the two
terminal couplings agree.  Since
\[
 m_K^2\epsilon_K
 =L^{(\upsilon-2z)K+O(1)}\longrightarrow0,
\]
its partition-function estimates are applicable for large $K$.
The per-volume scalars cancel separately in each rectangular test,
giving
\begin{align*}
 &\left|\Delta_{\beta}(u_1m_KL^N,u_2m_KL^N)
       -\Delta_{\widetilde\beta_K}(u_1m_KL^K,u_2m_KL^K)\right|\\
 &\hspace{35mm}\le C_Hm_K^2\epsilon_K\longrightarrow0
\end{align*}
uniformly over the displayed family.  Choose one sufficiently large
$K_0$ and set $M_0=m_{K_0}$.  Nonnegativity of the tests follows
from the two-direction transfer identities
\eqref{eq:trace-rectangle-identities}.  This proves
\eqref{eq:offset-box-test}.
\end{proof}

\subsection{Transfer bounds and the plane limit}

Only the finite list of tests in Lemma~\ref{lem:offset-box-test}
is needed from renormalization for the next result.  The rest is the
positive-transfer argument of Section~\ref{sec:trace}, now applied
uniformly to the offset family.

\begin{proposition}[Uniform volume comparison and slab decay]
\label{prop:uniform-volume}
Fix $K_0,M_0$ as in Lemma~\ref{lem:offset-box-test}.  There exist
$C,c>0$, independent of $N\ge K_0$, $|s|\le R$, and $k\ge0$,
such that the following hold for $\beta=\vartheta_N+s$.

Let $(u_1,u_2)\in\mathcal A$ and
$(n,w)=(u_1,u_2)2^kM_0L^N$.  If a bounded local observable $F$
has support contained in a rectangle of time extent at most $n/2$
and spatial extent at most $w/2$, then
\begin{equation}
 |\mu_{\beta;n,w}(F)-\mu_\beta(F)|
       \le C\|F\|_\infty e^{-c2^k}.
 \label{eq:offset-uniform-volume}
\end{equation}
The same bound holds when time is closed by any integer spatial
translation, provided the support has a lift satisfying these two
extent bounds.  The constant is uniform in the closing translation.

If two bounded local plane observables $F,G$ have supporting slabs
separated by $d\ge0$ transfer edges, in either microscopic axial
direction, then
\begin{equation}
 |\Cov_{\mu_\beta}(F,G)|
 \le \|F\|_\infty\|G\|_\infty
           e^{-(\log2)d/(M_0L^N)}.
 \label{eq:offset-slab-decay}
\end{equation}
For complex observables the covariance uses $\overline F G$.
\end{proposition}

\begin{proof}
For clarity we specify which parts of the transfer proof are
independent of the chosen trajectory.  The nearest-neighbor operator
$K_w$ of \eqref{eq:trace-kernel} is positive semidefinite and has
strictly positive kernel for every $\beta\ge0$.  Its largest
eigenvalue is simple; let $T_w=K_w/\lambda_1(w)$, let $P_w$ be the
projection onto its normalized positive ground state $\Omega_w$,
and write
\[
 s_j(w)=\operatorname{Tr}(T_w^j-P_w),\qquad j\ge1.
\]
Lemma~\ref{lem:trace-concentration} and
Proposition~\ref{prop:trace-doubling} use only this positivity and
the rectangular test.  They give
\[
 \Delta_\beta(2^kn,2^kw)\le64^{-1}2^{-2^k}
\]
from each test in \eqref{eq:offset-box-test}.  Applied at the
auxiliary aspects of \eqref{eq:offset-aspect-inventory}, this implies
\[
 s_{n/2}(w)\le Ce^{-c2^k},
 \qquad
 s_{w/2}(2n)\le Ce^{-c2^k},
\]
where the second expression refers to transfer in the spatial
direction.

An insertion on a slab of $r\ge1$ steps has normalized operator
$A_F$ whose kernel is bounded in absolute value by
$\|F\|_\infty K_w^r$ before normalization.  Thus
$|A_Ff|\le\|F\|_\infty T_w^r|f|$ and
$\|A_F\|\le\|F\|_\infty$, also for complex $F$.
For $r=0$, $A_F$ is multiplication by $F$ and obeys the same norm bound.
For $r\le n/2$, separating the ground-state projection in the
transfer trace gives \eqref{eq:trace-torus-cylinder}:
\[
 |\mu_{\beta;n,w}(F)-\langle\Omega_w,A_F\Omega_w\rangle|
       \le2\|F\|_\infty s_{n/2}(w).
\]
Use this first to compare $(n,w)$ with $(2n,w)$, and then, after
interchanging the coordinates, to compare $(2n,w)$ with $(2n,2w)$.
The two preceding excited-trace estimates give an error
$C\|F\|_\infty e^{-c2^k}$.  Summing over all further simultaneous
doublings proves \eqref{eq:offset-uniform-volume}; the limit is
$\mu_\beta$ by Proposition~\ref{prop:preliminary}.  Bounded
measurable insertions are covered as in the proof of
Theorem~\ref{thm:trace-uniform}: finite-support marginals have a
uniform density bound at this fixed $\beta$, and approximation in
product area measure extends the identification of the local limit.

For a closing shift $h$, its unitary $U_h$ commutes with $T_w$ and
fixes $\Omega_w$.  Place the seam outside the supporting slab.
The denominator and numerator are respectively
\[
 \operatorname{Tr}(T_w^nU_h),\qquad
 \operatorname{Tr}(A_FT_w^{n-r}U_h).
\]
Their remainders after the ground-state contribution have absolute
values at most $s_n(w)$ and $\|F\|_\infty s_{n-r}(w)$, uniformly
in $h$.  The reference tests ensure $s_n(w)<1/2$.  Division thus
changes the cylinder expectation by at most
$4\|F\|_\infty s_{n/2}(w)$, proving the shifted assertion.

Finally set $n_0=M_0L^N$.  The square test gives, at all
circumferences $w=2^kn_0$,
\[
 \|T_w^d-P_w\|\le e^{-(\log2)d/n_0}
\]
by \eqref{eq:trace-width-gap}.  For ordered slab insertions the
cylinder covariance is
\[
 \langle\Omega_w,
       A_{\overline F}(T_w^d-P_w)A_G\Omega_w\rangle.
\]
The insertion norm bounds prove \eqref{eq:offset-slab-decay} on
these cylinders.  Their local limit is the periodic plane law by
the half-period trace estimates and the torus comparison already
proved.  Passing to that limit proves the plane estimate.
Interchanging coordinates gives the other axial direction.
\end{proof}

\subsection{Common physical tori for the inclined schemes}

The shifted closing condition has a concrete role: it puts the
ordinary and inclined embeddings on the same physical square tori.
We spell out the integer geometry to fix both the length factor and
the finite list of aspects used above.

\begin{lemma}[Inclined periods]
\label{lem:inclined-periods}
Let $M_k=2^kM_0$.  Embed an ordinary depth-$N$ run with spacing
$L^{-N}$ and axes the identity, and a type $\pm$ run with spacing
$a=L^{-N}/5$ and axes $O_\pm^{-1}$.  The square physical torus of
side $M_k$ has ordinary microscopic periods $M_kL^N e_1$ and
$M_kL^N e_2$.  For type $+$ its microscopic periods, in
$(\mathrm{space},\mathrm{time})$ coordinates, are
\begin{equation}
 v_1=M_kL^N(3,4),\qquad v_2=M_kL^N(-4,3).
 \label{eq:offset-inclined-periods}
\end{equation}
They are equivalent to spatial circumference $25M_kL^N$ and time
height $M_kL^N$, closed by the spatial shift $7M_kL^N$.
The reflected type has the same circumference and height, with
shift $-7M_kL^N$.  These periods are compatible with the initial
inclined blocking and the subsequent ordinary blockings, and have
terminal square periods $M_k$.

Let $F_N$ be any bounded observable whose physical support in the
chosen embedding is contained in one fixed compact set, and put
$\beta=\vartheta_N+s$, $|s|\le R$.  Its support has lifts obeying
the half-extent hypotheses of Proposition~\ref{prop:uniform-volume}
for all sufficiently large $k$, uniformly in $N\ge K_0$ and $s$.
Its torus expectation therefore differs from its plane expectation
by $C\|F_N\|_\infty e^{-c2^k}$ in all three embeddings.
\end{lemma}

\begin{proof}
The columns of $O_+$ are $(3,4)/5$ and $(-4,3)/5$.
Multiplication of \eqref{eq:offset-inclined-periods} by
$(L^{-N}/5)O_+^{-1}$ gives $M_ke_1,M_ke_2$.
The integer basis change
\[
 3v_1-4v_2=(25M_kL^N,0),\qquad
 v_1-v_2=(7M_kL^N,M_kL^N)
\]
has determinant one.  It therefore preserves the period lattice
and proves the asserted shifted representation.  Reflection changes
the sign of the spatial shift and gives the type $-$ statement.

The initial coarse translation lattice of the inclined observation
is $5L O_\pm\Z^2$, as in
Lemma~\ref{lem:free-inclined-geometry}.  The periods above belong
to it and become $M_kL^{N-1}e_1,M_kL^{N-1}e_2$ in the output
frame.  The remaining ordinary steps yield square periods $M_k$.
The fixed divisibilities of $M_0$ and the terminal lower bound
ensure admissibility at every layer, by
Section~\ref{setup:scales}.

A compact physical support of diameter $D$ has microscopic
coordinate extents at most $C D/a\le C'DL^N$ in either inclined
frame; translations of its chosen lift do not change those extents.
For all sufficiently large $k$, these are at most half the time
height $M_kL^N$ and half the spatial circumference $25M_kL^N$.
The same assertion for the ordinary frame uses aspect $(1,1)$.
Thus the two aspects in $\mathcal A$ and the shifted estimate of
Proposition~\ref{prop:uniform-volume} prove the final claim.
\end{proof}
